% Transcription par lecture visuelle directe, sans OCR.
% Statut : transcrit ; relecture indépendante à effectuer.
% Source : PDF original, vue 093, côté gauche, page manuscrite 186, après le chapitre sixième.
\chapter{Préfecture de M. Philpin}
\begin{center}
Suspension de la célébration des fêtes dans une chapelle particulière, \&…
\end{center}

\section*{Sa nomination}
% Mention marginale de la page manuscrite 186.
La Congrégation conserva M. Chaulet à sa tête aussi longtemps que son règlement lui en laissa la faculté. Son administration avait été couronnée du plus heureux succès ; on avait pu considérer le temps de sa préfecture comme une époque de prospérité pour notre société. Elle s'était augmentée et fortifiée, avait mis la dernière main à son règlement ; ses sections s'étaient lancées dans la carrière des bonnes œuvres. Au dehors des alliances avantageuses avaient été contractées ; enfin la Congrégation avait signalé son influence d'une manière bien heureuse par la part qu'elle avait prise à la formation des confréries du St Sacrement et aux honneurs rendus au souverain pontife. Tous les membres de la Congrégation
% Source : PDF original, vue 093, côté droit, page manuscrite 187.
eussent donc désiré pouvoir conserver encore le préfet sous le gouvernement duquel tant de bien s'était opéré. Mais le règlement était positif ; à la fin de la troisième année d'exercice, le préfet ne pouvait plus être réélu ; au 15 août 1806, époque de l'élection, il fallut donc songer à lui donner un successeur. Le choix de la Congrégation se fixa sur M. Philpin, qui était alors secrétaire. Comme cet ami que l'on appelle aujourd'hui M. Philpin de Rivière, n'habite plus Lyon depuis longtemps et quoiqu'il reparoisse encore au milieu de nous lorsque ses affaires l'appellent dans notre ville, il pourrait se faire que plusieurs d'entre vous n'en eussent pas même entendu parler, je dois vous faire entrer en connaissance avec lui, en vous disant ici quelques mots de ce qui le concerne.

M. Philpin n'appartenait point par sa naissance à notre ville. Ayant émigré au commencement de la Révolution avec sa famille, il était rentré en France dans un moment où on pouvait espérer le retour de l'ordre. Il vint à Lyon pour se dérober plus facilement à l'œil inquisiteur de la police révolutionnaire. Sous le nom d'Antoine qui était son nom de baptême, il s'était associé avec un excellent chrétien et trouvait dans le commerce des moyens d'existence. Malgré sa modestie et l'obscurité dans laquelle il s'enveloppait, les vertus de M. Philpin ne tardèrent pas à le faire apprécier par tous ceux qui furent à portée d'en être les témoins. Lorsque les confréries du St Sacrement se formèrent, il fut un des premiers appelé à faire partie de celle de St Pierre. Sa piété s'y distingua ; il y fut nommé maître des cérémonies, place qu'il a remplie avec zèle et distinction. Cette circonstance le mit en rapports fort intimes avec ceux des congréganistes qui faisaient avec lui partie du bureau de la confrérie. Il n'était pas difficile
% Source : PDF original, vue 094, côté gauche, page manuscrite 188.
de juger combien un semblable sujet pouvait être utile à la Congrégation ; il y fut donc proposé et admis. Procurer à M. Philpin un moyen de plus de travailler à la gloire de Dieu c'était très certainement satisfaire le plus vif désir de son cœur, aussi s'empressa-t-il d'accepter, et la Congrégation fut alors à portée de connaître tout le prix de l'excellente acquisition qu'elle venait de faire. Piété, zèle, capacité, douceur, modestie et fermeté, tout ce qui peut faire à un degré éminent un bon congréganiste, se trouvait réuni dans cet excellent ami.

Permettez-moi de vous citer de lui une seule parole ; elle le peindra mieux que tout ce que je pourrais ajouter. Il fut placé à son entrée dans la Congrégation, dans la section des prisons et hôpitaux ; j'en faisais alors partie ; je le vis se porter avec tant d'ardeur aux œuvres de cette section, surtout en ce qui concernait le soulagement des prisonniers, que je ne mis pas en doute qu'il avait un goût particulier pour cette œuvre, et qu'elle lui était extrêmement agréable. Un jour je lui exprimai dans une conversation intime combien je me réjouissais de voir que nos chefs avaient aussi bien deviné son inclination, et l'avaient ainsi placé tout à fait à sa convenance. Il mit dans sa réponse un certain abandon. « Oui, me répliqua-t-il, je me suis estimé très heureux d'avoir été chargé de cette œuvre ; car certainement de toutes les œuvres de miséricorde c'est celle pour laquelle j'éprouve le plus de répugnance ; je suis donc bien sûr alors que l'esprit propre n'y sera pour rien, et qu'en m'efforçant de vaincre dans des vues de foi et d'obéissance les dégoûts qu'elle m'inspire, j'en serai plus agréable au bon Dieu qui a manifesté sa volonté par celle de nos supérieurs. » C'est cet homme qui savait si bien obéir dans le seul désir de plaire à notre bon
% Source : PDF original, vue 094, côté droit, page manuscrite 189.
maître, que le Seigneur appelait à diriger la marche de la Congrégation. Il fut installé comme à l'ordinaire, le jour de la fête de la Nativité de la Ste Vierge.

\section*{Suspension des réunions pour les fêtes}
% Mention marginale de la page manuscrite 189.
Tout jusqu'ici avait souri à la Congrégation ; mais il était facile de s'apercevoir que bientôt les épreuves allaient commencer. Le gouvernement s'apprêtait à devenir prochainement persécuteur, en attendant il se montrait ombrageux. Cette disposition de sa part se fit bientôt sentir à la Congrégation. M. Philpin était à peine entré en fonctions, qu'un premier sacrifice fut exigé de nous, ce fut celui de nos fêtes, de ces belles solennités dans lesquelles nos cœurs éprouvaient tant de jouissances. Quoique l'existence de notre société fût demeurée secrète, néanmoins ce qui malgré tous nos efforts transpirait dans le public, commençait à éveiller l'attention. Si d'un côté ces douces récréations consolaient et édifiaient les âmes fidèles, de l'autre elles attiraient les regards d'une police inquisitoriale que tout inquiétait. Le cardinal craignit que nos réunions particulières dans une chapelle où les initiés seuls étaient admis ne parussent suspectes à cette police et il se résigna à les sacrifier. Notre bon père Roger fut chargé de nous signifier cet arrêt ; il fallut bien se soumettre. Il fallut jusqu'à un temps meilleur renoncer à cette chère chapelle des Chartreux, où depuis le premier anniversaire de la fondation, nos fêtes s'étaient célébrées avec tant de pompe, surtout avec tant de bonheur.

\section*{Les saints du mois}
% Mention marginale de la page manuscrite 190 ; reportée au début du passage.
La Congrégation fut vivement affectée de cette privation ; mais elle ne perdit pas courage ; elle s'efforça de suppléer par tout ce qui dépendait d'elle aux jouissances qui lui étaient enlevées. On redoubla de régularité et d'exactitude dans la tenue des assemblées générales. Ce fut vers ce temps-là qu'on adopta l'usage de distribuer
% Source : PDF original, vue 095, côté gauche, page manuscrite 190.
les saints du mois à la fin de chaque assemblée ordinaire. Vous savez que nous appelons les saints du mois ces petites légendes accompagnées de l'image du saint qui en est l'objet, d'une méditation et de l'indication d'une vertu, qu'on continue à distribuer à la fin de nos séances. Cette distribution est destinée à remplacer pour les congréganistes, les jetons de présence qu'on est en usage de distribuer dans certaines compagnies ou administrations comme une prime accordée à l'exactitude ; mais en établissant cette pieuse pratique, la Congrégation avait en vue un autre but que je crains qu'il ne soit aujourd'hui un peu oublié ; il est donc bon de le rappeler. La Congrégation s'était proposé de former entre tous ses membres une sorte d'union perpétuelle par la prière et par la communion ; c'est pourquoi chaque congréganiste était invité en prenant pour patron du mois, le saint qui lui était échu, de communier en union avec la Congrégation, si cela lui était possible au jour même de la fête du saint ou du moins le dimanche avant ou le dimanche après. Jusqu'à l'assemblée suivante, les congréganistes étaient également invités à prendre pour sujet de leurs méditations particulières soit celle qui leur était indiquée, soit la vie du saint soit la vertu dont la pratique leur était spécialement recommandée. Il est à présumer que les membres se conformèrent avec soin, dans la nouveauté de l'institution, aux avis qui leur furent donnés à ce sujet. Ce qui pourrait le faire croire, c'est l'intérêt qu'ils mettaient tous à ne pas manquer de recevoir leurs saints du mois. La crainte d'en être privé eut l'heureux avantage de les rendre beaucoup plus exacts de se trouver aux assemblées.

\section*{Mode de célébrer les fêtes}
% Mention marginale des pages manuscrites 191–192.
Quoique la Congrégation fût privée du bonheur de solenniser
% Source : PDF original, vue 095, côté droit, page manuscrite 191.
ses fêtes avec l'éclat dont jusqu'alors elle les avait entourées, elle ne renonça pas pour cela à les célébrer ; elle ne renonça pas non plus à faire des réceptions ; elle se vit donc obligée de recourir à des formes nouvelles pour remplacer celles dont l'usage lui était interdit. Voici les coutumes qui furent établies à cette époque et dans lesquelles on a persévéré jusqu'au moment où il a été possible de revenir à l'ancienne méthode. À l'approche de chaque fête on indiquait l'heure de la communion générale, et l'église où on se réunirait pour cet objet. Le plus souvent c'était Fourvières. Nous nous trouvions heureux lorsque nous étions rassemblés dans ce sanctuaire de Marie. Si les voix ne pouvaient plus confondre leurs accents pour chanter d'un commun accord les miséricordes de notre Dieu et les grandeurs de notre mère, du moins les cœurs plus unis que jamais renouvellaient en commun et leur amour et leur reconnaissance. Là les foibles étaient encouragés et fortifiés, car ils se sentaient soutenus par les prières de leurs frères, dont une bien douce association leur faisait partager tous les mérites. Avant la communion ceux qui avaient été jugés dignes d'entrer dans la famille de Marie, prononçaient à voix basse en présence des assistants, qui avaient toujours le soin de se rapprocher d'eux, le renouvellement des promesses de leur baptême et leur consécration à la Ste Vierge. La sainte eucharistie, ce pain de l'union de tous les fidèles avec leur Dieu, venait par sa divine présence consacrer l'alliance qu'ils contractaient avec leurs nouveaux frères.

Dans la soirée il y avait une assemblée extraordinaire ; ces assemblées de fête étaient accompagnées d'un peu de solennité. Les membres du bureau s'y rendaient toujours vêtus de noir (c'était alors la grande tenue), l'image de la Ste Vierge paraissait au milieu de l'assemblée comme le point de centre qui nous unissait tous. Elle était placée
% Source : PDF original, vue 096, côté gauche, page manuscrite 192.
sur une crédence décorée avec soin et environnée de flambeaux et de fleurs ; à ses côtés on faisait placer les nouveaux reçus ; comme les plus jeunes de la famille, ils avaient le privilège d'être sous les ailes de la mère commune.

Le préfet ouvrait la séance par la récitation du Veni Creator ; ensuite au lieu du procès-verbal ordinaire, le secrétaire allait chercher dans un vieux registre celui de la même fête qui nous réunissait, célébrée dans une des précédentes années avec tout l'appareil que nous pouvions y mettre lorsque nous jouissions d'une entière liberté. C'est ainsi qu'on cherchait à se consoler du manque présent par le souvenir des jouissances du passé. On récitait ensuite vêpres et complies de l'office de la Ste Vierge. Le préfet adressait une allocution aux nouveaux reçus, après laquelle toute l'assemblée se tenant debout ils venaient successivement prononcer de nouveau aux pieds de Marie à haute voix les mêmes actes de renouvellement des vœux du baptême et de consécration à la Ste Vierge que déjà ils avaient faits le matin à voix basse aux pieds des autels. Tous leurs nouveaux frères les serraient ensuite tour à tour entre leurs bras en répétant de cœur plus encore que de bouche ces paroles du psalmiste consacrées par l'usage à cette cérémonie : \textit{Ecce quam bonum et quam jucundum habitare fratres in unum}\footnote{(Note du transcripteur) Littéralement : « Voici comme il est bon et comme il est agréable que les frères habitent ensemble ». Psaume 132 (133), verset 1.} ; enfin le cantique de Marie, le Magnificat terminait et couronnait la fête.

On employait à peu près les mêmes formes pour toutes les cérémonies qui auraient dû suivant nos anciens usages se passer dans la chapelle, comme la lecture du bref d'indulgence au 4\up{e} dimanche après Pâques, l'installation du préfet au jour de la Nativité ; toutes les prières qui dans ces cas-là auraient été chantées dans la chapelle étaient récitées en assemblée générale.

% Source : PDF original, vue 096, côté droit, page manuscrite 193.
\section*{Visite de Mgr Fournier, évêque de Montpellier}
% Mention marginale de la page manuscrite 194 ; reportée au début du récit.
Quelques circonstances particulières pouvaient aussi parfois donner occasion à la convocation d'assemblées extraordinaires comme cela arriva, par exemple, à l'époque du passage de M. Fournier, ce célèbre orateur dont nous avons déjà parlé qui venait d'être nommé évêque de Montpellier.

La liberté toute évangélique avec laquelle cet éloquent orateur annonçait la parole de Dieu déplaisait très fort au despote qui gouvernait la France. Toutefois sa tyrannie n'avait point encore atteint le degré auquel elle est parvenue plus tard. M. Fournier savait se renfermer assez dans les bornes de la prudence pour qu'il n'offrît aucun prétexte plausible de sévir contre lui. Au lieu d'user de violence, Napoléon employa un moyen plus doux pour le réduire au silence, ou du moins pour circonscrire la puissance de sa parole dans des bornes étroites ; ce moyen fut de l'appeler à l'épiscopat ; il le nomma évêque de Montpellier et effectivement depuis cette époque son imposante voix n'a plus fait retentir d'autres échos que ceux de sa cathédrale.

Le nouvel évêque passa à Lyon en se rendant à Montpellier pour aller prendre possession de son siège. La Congrégation se rappelait avec une vive reconnaissance la bonté extrême que ce prélat lui avait témoignée dans son enfance en ne dédaignant pas de faire entendre sa voix à une société qui avait alors si peu de droits à une telle faveur. Nos autorités crurent donc qu'il était de leur devoir de lui faire pendant le séjour qu'il fit à Lyon une visite comme un témoignage de notre gratitude. Il les reçut avec tant de bonté que cela les encouragea à oser prier sa grandeur d'honorer encore la Congrégation de sa présence. Le prélat accueillit cette demande très favorablement, il fit plus ; il promit de prononcer en
% Source : PDF original, vue 097, côté gauche, page manuscrite 194.
milieu de nous un discours inédit qu'il ne jugeait pas convenable de donner en chaire. On prit donc jour et heure, et on s'empressa de convoquer une assemblée qu'on entoura de toute la solennité qu'on put lui donner. C'était la première fois que la Congrégation allait être honorée de la présence d'un évêque. Elle voulut mettre à sa réception tout l'appareil qui dépendait d'elle. Voici le cérémonial qui fut adopté et qui fut exactement suivi ; tous les membres furent avertis de se rendre à l'assemblée en grande tenue ; quoiqu'au milieu du jour l'assemblée se tint aux flambeaux, outre ceux qui brûlaient autour de la Ste Vierge des candélabres garnis de bougies étaient placés de distance en distance et notamment sur la table du bureau. Deux maîtres des cérémonies furent délégués pour aller chercher Monseigneur à l'archevêché où il logeait et l'amener en voiture. Quatre congréganistes en habits noirs et en gants blancs le reçurent au bas de l'escalier de la maison où devait se tenir l'assemblée. Le père, le préfet et les autres membres du bureau le reçurent à leur tour à la porte de l'appartement, l'introduisirent dans la salle d'assemblée et le conduisirent au fauteuil qui lui était destiné. Il était placé au centre de l'assemblée à la place qui est ordinairement occupée par le préfet. Le préfet se plaça à la droite de Mgr et le père à sa gauche ; c'est ainsi que cela se pratique encore toutes les fois que nous recevons un évêque. Le prélat fut prié de faire les prières d'ouverture et de clôture. Cette séance dura environ trois heures. Mgr nous donna lecture d'un très beau discours dans lequel il démontra avec évidence et avec tout le charme de sa brillante élocution, l'accord parfait qui existe entre les connaissances acquises par la physique ou la géologie,
% Source : PDF original, vue 097, côté droit, page manuscrite 195.
et le grand événement du déluge tel que nous en lisons l'histoire dans les livres saints. Mgr fut raccompagné avec le même cérémonial avec lequel il avait été reçu ; il nous laissa tous pénétrés de sa bonté et remplis de joie d'avoir reçu une semblable visite dont le souvenir doit être conservé dans les fastes de la Congrégation.

Nos assemblées se tenaient alternativement chez ceux de nos amis qui se trouvaient le mieux en position de les recevoir. Ceux chez lesquels nous nous réunissions le plus souvent à cette époque étaient MM. Ducoin et Brette. Nous avons déjà parlé de M. Ducoin, bientôt nous en reparlerons encore, mais puisque je viens de rappeler le nom de M. Brette, permettez-moi de vous entretenir quelques instants de ce digne ami, dont la sainteté a été pendant longtemps un grand sujet d'édification pour la Congrégation.

\section*{M. Brette}
% Mention marginale de la page manuscrite 195 ; placée au début du portrait.
Le docteur Brette était né à Grenoble ; il s'était consacré tout jeune à l'art de la médecine ; élève interne à l'Hôtel-Dieu de Lyon il s'était fait respecter et admirer de tous ses confrères dont un grand nombre cependant se trouvait malheureusement bien éloigné de partager ses sentiments religieux. Il fut ensuite à Paris pour compléter ses études en suivant un cours de médecine ; il fut reçu dans la Congrégation de Paris pendant son séjour dans cette capitale. Cette circonstance aurait suffi pour lui donner ses entrées dans les assemblées de notre Congrégation lorsqu'il venait se fixer à Lyon, mais ce n'était pas assez pour nous. Sa réputation de piété, on pourrait presque dire, sa réputation de sainteté l'avait précédé ; nous étions jaloux de voir l'intimité la plus entière et la plus complète devenir le lien de l'union que nous formerions avec lui ; aussi la Congrégation s'empressa-t-elle de le recevoir, et ce ne fut pas comme congréganiste de Paris, mais
% Source : PDF original, vue 098, côté gauche, page manuscrite 196.
comme congréganiste de Lyon qu'il fut admis.
La Providence n'avait pas destiné M. Brette à diriger la Congrégation, ou à y occuper les emplois ; il était appelé à lui rendre de bien plus éminents services en la soutenant par la ferveur de ses prières, par l'esprit intérieur dont il était animé, et par toutes les vertus dont il mettait sous les yeux de tous les congréganistes un si parfait modèle. On ne peut pas mettre en doute que son humilité et son admirable simplicité ne fussent d'un bien plus grand prix aux yeux de Dieu et ne contribuassent bien plus à mériter à la Congrégation la protection du ciel, que les saillies du zèle ardent de plusieurs d'entre nous, où l'amour-propre ne vient que trop souvent mêler un funeste alliage. Je ne sais si je me trompe, mais j'ai toujours eu la tentation de comparer M. Brette à cet heureux disciple qui mérita d'entendre cet éloge sorti de la bouche même de notre Seigneur : celui-ci est un véritable Israélite dans lequel il n'y a point de malice. Toujours pénétré de l'esprit de foi et d'une vive confiance en Dieu, on le voyait souvent au pied du lit de ses malades lorsque les ressources qu'une science humaine pouvait lui offrir pour leur rendre la vie et la santé lui paraissaient épuisées, fléchir les genoux et implorer le secours du seul médecin véritable. Souvent alors il arrivait que Dieu accordait à sa foi et à l'ardeur de ses prières ce qu'il avait refusé à l'efficacité des remèdes. N'en doutons pas, Messieurs, ce sont les hommes qui ressemblent à M. Brette qui font la véritable force et sont les plus fermes appuis d'une société comme la nôtre.

\section*{Mgr de Forbin-Janson}
% Mention marginale de la page manuscrite 197 ; reportée au début du passage.
La Congrégation, malgré les contrariétés qu'elle éprouvait, grâce à la piété éminente de quelques-uns de ses membres et au
% Source : PDF original, vue 098, côté droit, page manuscrite 197.
zèle de la plupart d'entre eux continuait suivant sa devise à marcher toujours en avant. Les sections s'assemblaient avec exactitude, elles s'efforçaient d'accroître de plus en plus la somme de leurs bonnes œuvres. Lorsque les congréganistes étrangers qui nous étaient associés, faisaient quelque séjour dans notre ville, ils sollicitaient et obtenaient précisément la faveur de venir seconder et animer le zèle lyonnais, par l'exemple du leur. Parmi ceux-ci nous distinguerons surtout M. de Forbin-Janson, membre de la Congrégation de Paris et qu'en cette qualité nous avions associé à la nôtre. Dans le temps dont nous nous occupons, il fit à Lyon un séjour fort long, je crois de près de deux ans. Il demanda avec instance et on lui accorda la permission de partager les travaux de la section des aumônes. Il mit à l'exercice de ces bonnes œuvres toute la vivacité de son caractère et tout le feu dévorant de son zèle. Vous savez, Messieurs, que plus tard, revêtu du sacerdoce, M. de Janson est devenu d'abord missionnaire et ensuite évêque. En cette double qualité toujours dévoré pour la cause de l'Évangile il s'est montré un véritable apôtre de J. C. Plusieurs fois, et encore dans ces dernières années, nous avons eu le bonheur de le voir reparaître dans nos assemblées et chaque fois il s'efforce de faire pénétrer dans nos cœurs cet amour de Dieu qui remplit le sien.

\section*{Établissement d'une société de bonnes œuvres à Paris}
% Mention marginale de la page manuscrite 198 ; reportée au début du passage.
M. de Noailles ne faisait pas à Lyon des séjours aussi prolongés, mais il y venait souvent et toujours heureux lorsqu'il retrouvait ses frères de Lyon et pouvait partager leurs bonnes œuvres, il ne se bornait pas à assister aux assemblées générales, il allait frapper à la porte des sections et on s'empressait de lui ouvrir. Personne n'admirait plus sincèrement que lui
% Source : PDF original, vue 099, côté gauche, page manuscrite 198.
le mécanisme de notre organisation intérieure, il eût bien désiré la faire goûter et admettre par la Congrégation de Paris ; il fit même à ce sujet plusieurs démarches, mais trop d'obstacles s'y opposaient : la manière d'être, le genre de vie de la plupart des congréganistes de Paris, l'esprit même de leur congrégation qui étaient tout à fait contraires. M. de Noailles se voyant dans l'impossibilité de réussir, chercha en désespoir de cause le moyen de former à côté de la Congrégation, une société de bonnes œuvres qui se diviserait en sections comme la nôtre et qui serait composée des congréganistes auxquels leur zèle inspirerait le désir d'en faire partie. Ce plan a eu un succès complet. Ainsi ce fut à l'imitation de la Congrégation de Lyon que prit naissance cet accessoire de la Congrégation de Paris. Cette institution a rendu de notables services. Une de ses sections surtout a été extrêmement utile, c'est celle qui est chargée du soin de cette masse d'enfants qui sous le nom de Savoyards affluent chaque année sur le pavé de Paris pour y chercher des moyens d'existence.

Le nombre des membres de la Congrégation augmentait graduellement quoique avec lenteur ; il se faisait pourtant beaucoup de réceptions, mais elles étaient presque compensées par les pertes successives des nombreux congréganistes que des circonstances diverses forçaient de quitter la ville. Les jeunes gens qui faisaient partie de la Congrégation n'étaient pas encore pour la plupart fixés définitivement par un état qui rendît le domicile où ils se trouveraient le lieu irrévocable de leur demeure. Le nombre annuel de ces émigrants peut s'évaluer à cette époque de dix à quinze ; vous concevez qu'il devait falloir bien des réceptions
% Source : PDF original, vue 099, côté droit, page manuscrite 199.
pour couvrir ce déficit ; cependant elles y suffisaient et au-delà ; sans compter le grand nombre de ses membres absents, la Congrégation avait alors de 40 à 50 membres qui prenaient part à ses œuvres. C'est ce que nous appellerions aujourd'hui membres en activité.

\section*{Négociation avec la Congrégation de la Ste Famille pour les transferts}
% Mention marginale de la page manuscrite 199.
Elle fut menacée, vers ce temps-là, par la force même des circonstances d'éprouver une diminution bien sensible. Sept à 8 de ses membres se trouvèrent mariés ; ce nouvel état qui leur imposait de nouvelles obligations, paraissait devoir les exclure de la Congrégation des jeunes gens, et marquer leur place dans celle de la Ste Famille composée tout exprès pour les hommes mariés. Cette société ne demandait pas mieux que d'ouvrir ses rangs à plusieurs d'entre eux, mais elle se réservait le droit de les soumettre à l'épreuve du scrutin et de refuser par conséquent ceux qu'elle ne jugerait pas à propos d'admettre. C'eût été exposer plusieurs de nos amis à une espèce d'ostracisme et les priver pour le reste de leur vie des avantages attachés à la qualité de congréganiste. Ceux donc à qui le scrutin de ces Messieurs fut favorable crurent ne pas devoir ainsi abandonner leurs frères. Ils déclarèrent qu'ils ne consentiraient à changer de Congrégation, qu'autant que tous leurs confrères mariés en changeraient avec eux, et que de plus la Congrégation des hommes prendrait pour l'avenir l'engagement de recevoir sans aucune formalité les congréganistes jeunes gens à l'époque de leur mariage. Cette proposition fut longuement discutée dans la Congrégation de la Ste Famille, mais elle n'y fut point acceptée ; par ce refus la question fut définitivement jugée ; ceux d'entre nous qui étaient mariés se décidèrent avec joie à rester avec les jeunes gens et même à y vieillir. Nous verrons plus tard comment la Providence qui a toujours veillé sur la Congrégation avec un soin si paternel a amené par la suite
% Source : PDF original, vue 100, côté gauche, page manuscrite 200.
une heureuse solution à cette difficulté qui de tout côté présentait des inconvénients.

\section*{Neuvaines à Fourvières}
% Mention marginale de la page manuscrite 200.
J'ai encore à vous parler ici de Fourvières, Messieurs, vous n'en serez pas étonnés, car vous devez comprendre que depuis que ce sanctuaire avait été rendu à nos désirs, les congréganistes enfants de Marie, aimaient se plaire à y venir invoquer Marie ; c'était un dédommagement que la Providence leur avait préparé, avant de les priver de la célébration de leurs fêtes. Aussi non seulement ils y allaient souvent prier isolément, non seulement, comme nous l'avons déjà dit, on y indiquait fréquemment des communions générales, mais encore dès que quelque inquiétude nouvelle venait agiter nos cœurs, on s'empressait de recourir à Marie ; le préfet indiquait alors une neuvaine à Fourvières et tous les membres étaient fidèles à ce pieux rendez-vous. On avait recours à ce moyen efficace, toutes les fois que le tirage au sort de la conscription venait nous menacer de nous enlever quelques-uns de nos frères. De même qu'à l'approche du loup dévorant qui tourne autour du troupeau, prêt à s'élancer sur sa proie, les timides agneaux se pressent autour de leur mère pour y chercher un refuge, de même alors les congréganistes entouraient l'autel de Marie pour la supplier de sauver ses enfants. Leur confiance n'a jamais été déçue ; une protection visible s'est constamment manifestée à cet égard sur la Congrégation, dans un temps où la cherté des remplaçants eût rendu ce moyen impossible à plusieurs d'entre nous. Tous, excepté un seul, sont parvenus d'une manière ou d'une autre à échapper au danger ; le seul que Dieu appelait sans doute, à l'état militaire est M. Michaudon, cadet. Si la Ste Vierge
% Source : PDF original, vue 100, côté droit, page manuscrite 201.
le laisse partir, du moins elle a veillé sur lui. Demeuré au service pendant de longues années, il est parvenu au grade de capitaine, dans un temps où l'ange exterminateur semblait avoir reçu mission de marcher à la suite de nos armées pour les décimer. Il est demeuré constamment sain et sauf. Il est vrai que M. Michaudon dans sa vie militaire, n'avait oublié ni Fourvières, ni la Congrégation. Lorsqu'il entrait en campagne, il ne manquait jamais de se faire recommander aux prières de ses amis, et au moment d'une bataille il considérait en esprit le clocher de Fourvières, il appelait Marie à son secours, il se plaçait sous les ailes de sa mère, et grâce à cette puissante protection qu'il n'a jamais méconnue, il a été préservé des périls de plus d'un genre dont il était environné.

\section*{Célébration du 3\up{e} dimanche de juillet}
% Mention marginale des pages manuscrites 201–202.
Pour compléter le tableau de la marche de la Congrégation depuis la suspension de nos fêtes, il me reste à vous dire comme se célébrait à cette époque le 3\up{e} dimanche de juillet. Le matin, comme dans les autres fêtes, on se bornait à une communion générale dans une église. C'était assez ordinairement à Saint-Bonaventure, où nous passions à l'ombre du concours que la fête patronale attirait. L'après-midi les congréganistes se réunissaient dans une maison de campagne du voisinage de la ville, c'était assez souvent à cette époque une maison du faubourg de Serin appelée la Venerie qui appartenait aux parents d'un de nos amis. Il y avait une chapelle et de fort agréables ombrages. On s'y rendait le cœur plein de la douce joie du bonheur que trouvent des frères dans une réunion que Dieu lui-même a formée et leur esprit disposé à donner un libre cours à tous les transports que les souvenirs que cette journée rappelle pouvaient exciter.
On commençait la cérémonie en récitant ensemble dans la chapelle
% Source : PDF original, vue 101, côté gauche, page manuscrite 202.
vêpres et complies de la Ste Vierge ; ensuite les litanies, en renouvelant la consécration, enfin on récitait le Te Deum. On passait ensuite sous les arbres et on ouvrait l'assemblée par les prières d'usage. Pendant tout le temps que durait cette assemblée, la Congrégation se tenait debout, formant un cercle. Le secrétaire ouvrait son registre et lisait le procès-verbal de la fondation ; il faisait suivre cette lecture de celle du procès-verbal d'un des précédents anniversaires. Le préfet ou un membre désigné par lui jetait ensuite au milieu de l'assemblée quelques paroles toutes imprégnées de ces sentiments de reconnaissance et d'amour qui agitaient tous les cœurs. On proclamait la liste des candidats pour l'élection du préfet ; s'il était arrivé quelques lettres des membres absents on en donnait lecture. On faisait la quête pour les pauvres, réservée comme aujourd'hui, à la section des aumônes. Le père ajoutait avec chaleur quelques mots toujours qui s'échappaient de son cœur comme des traits enflammés. Arrivait ensuite le baiser de charité que chacun donnait et recevait tour à tour. Cette séance était suivie d'un repas maintenu peut-être mieux qu'il ne l'est à présent, dans les limites de la modestie chrétienne, mais qui était tout aussi joyeux ; les chansons, les récits, les pieux élans de dévotion envers Marie, ceux de la tendre et fraternelle amitié qui unissait entre eux tous les congréganistes, terminaient cette heureuse journée. Elle était pour nous un reste précieux de nos solennités des premières années ; nous le savourions avec délices. Mais hélas ! l'épreuve qui nous était réservée était à peine commencée. Elle était encore loin d'être parvenue au point où elle devait atteindre.
